% Generated draft: requires mathematical and source review.
\documentclass[11pt]{article}
\usepackage[margin=25mm]{geometry}
\usepackage{fontspec}
\setmainfont{Latin Modern Roman}
\usepackage{amsmath,amssymb,mathtools,unicode-math}
\setmathfont{Latin Modern Math}
\usepackage{enumitem,needspace,xurl,hyperref}
\setlength{\parindent}{0pt}
\setlength{\parskip}{5pt}
\setlength{\emergencystretch}{3em}
\newcommand{\N}{\mathbb{N}}
\newcommand{\Z}{\mathbb{Z}}
\newcommand{\Q}{\mathbb{Q}}
\newcommand{\R}{\mathbb{R}}
\newcommand{\C}{\mathbb{C}}
\newcommand{\F}{\mathbb{F}}
\newcommand{\A}{\mathbb{A}}
\newcommand{\PP}{\mathbb{P}}
\newcommand{\E}{\mathbb{E}}
\newcommand{\eps}{\varepsilon}
\newcommand{\dd}{\,\mathrm{d}}
\begin{document}
\section*{UnsolvedMath ID: 8700039}
\textbf{Source problem code:} catalogue entry
\subsection{Definitions and mathematical statement}
Let \(\mathcal E_1,\ldots,\mathcal E_n\) be pairwise nonisogenous elliptic curves over \(\C\). For each \(h\), choose fundamental periods \(\omega_{1h},\omega_{2h}\) for the associated Weierstrass function \(\wp_h\), and associated quasiperiods \(\eta_{1h},\eta_{2h}\). For example, in the usual normalization the Weierstrass zeta function satisfies \(\zeta'_h=-\wp_h\) and \(\zeta_h(z+\omega_{jh})-\zeta_h(z)=\eta_{jh}\). Let \(P_{ih}\in\mathcal E_h(\C)\), \(1\leq i\leq r_h\), and let \(p_{ih}\) and \(d_{ih}\) be their associated elliptic integrals of the first and second kinds, respectively, as specified by the source's missing normalization conventions.

The statement uses a field \(k_h\) without defining it. The conventional intended choice is \(k_h=\operatorname{End}(\mathcal E_h)\otimes_\Z\Q\), embedded in \(\C\) by its action on the analytic uniformization; this interpretation requires confirmation. Put \(\kappa_h=[k_h:\Q]\), and let
\[
 d_h=\dim_{k_h}\operatorname{span}_{k_h}
 \{\overline p_{1h},\ldots,\overline p_{r_h,h}\}
 \quad\text{in }\C/(k_h\omega_{1h}+k_h\omega_{2h}).
\]
Here the bars denote classes in that quotient. The integer \(d_h\) is distinct from the second-kind integrals \(d_{ih}\).

Bertolin's conjecture as recorded here asserts
\[
 \operatorname{trdeg}_{\Q}\Q\Bigl(
 \{j(\mathcal E_h),\omega_{1h},\omega_{2h},\eta_{1h},\eta_{2h}\}_{h=1}^n,
 \{P_{ih},p_{ih},d_{ih}\}_{1\leq h\leq n,\ 1\leq i\leq r_h}
 \Bigr)
 \ \geq\ 2\sum_{h=1}^n d_h+4\sum_{h=1}^n\kappa_h^{-1}-n+1.
\]
The notation adjoining \(P_{ih}\) must mean adjoining coordinates in specified models, since a point itself is not a complex field element; the models and treatment of the identity point are not supplied. No algebraicity of the curves, their invariants, or their chosen points is imposed by the excerpt. Nor is a normalization of the second-kind integrals given. These omissions must be resolved without reducing the conjecture to a statement only about periods.
\subsection{Short English statement}
For pairwise nonisogenous complex elliptic curves and chosen points, Bertolin predicts a specific transcendence-degree lower bound involving their invariants, periods, quasiperiods, point coordinates, and first- and second-kind elliptic integrals. The endomorphism-field and coordinate conventions need confirmation.
\subsection{Sources}
\begin{itemize}
\item Local UnsolvedMath registry, permanent ID 8700039, source code catalogue entry; source title ``Conjecture 3.21 --- Bertolin''; registry attribution Michel Waldschmidt; no proposal year supplied; updated 2026-08-06.
\item Waldschmidt, \emph{Open Diophantine Problems} (2004), Conjecture 3.21: \url{https://arxiv.org/abs/math/0312440}. Registry access date 2026-07-29; source-TeX extraction. Current primary-literature status and redistribution terms remain marked for review.
\item Bertolin is the supplied conjecture attribution. The summary mentions broader period and motivic conjectures and weak partial results without supplying an exact further reference or verified resolution.
\end{itemize}
\end{document}
